
\section{Twenty-First Sunday after Trinity: Always to pray and not grow weary}

\begin{quote}

Luke 18:1-8. And He spoke to them also a parable to this end, that one ought always to pray and not grow weary. And He said: There was a judge in a certain city, who did not fear God and had no regard for any man. And there was a widow in that same city, and she came to him and said: Vindicate me against my adversary. And for a long time he would not; but afterward he said within himself: Though I do not fear God nor regard any man, yet because this widow causes me much trouble, I will vindicate her, lest by her continual coming she wear me out. And the Lord said: Hear what the unjust judge says! And shall not God vindicate His elect, who cry to Him day and night, though He delays long over them? I tell you: He shall speedily vindicate them. Nevertheless, when the Son of Man comes, will He find faith on the earth?

\end{quote}

\bigskip

A widow suffering wrong is the image that the Lord chooses in order to set forth the believing soul and the believing Church on earth. And this widow prayed and begged and begged and prayed the ungodly and unjust judge to vindicate her, until she obtained what she desired. It was not her good cause that helped her; it was her tireless persistence in prayer. The judge cared nothing for her good cause; but her stubborn endurance became to him vexation and trouble, and so he helped her.

It is not the Savior’s meaning hereby to show how God is toward us; for God is precisely not like an unjust judge. Rather, the meaning is to show us how we are to be steadfast in prayer. For if persistent prayer can overcome the perverse unwillingness of an unjust judge, how much more shall persistent prayer overcome the heart of a loving Father and obtain the hearing of prayer.

A widow with an evil and unjust adversary is a moving image of the believing soul in the world. It is often truly lonely and forsaken; for it has bidden the world farewell, and it is dead to the world as the world is dead to it. It is homeless in the world, and as yet only on the way to the heavenly home. It has not yet reached home; it yearns and longs for it. And while it is thus in widowhood and forsaken, it has an evil adversary, who always tries to place himself between it and its heavenly inheritance, an accuser who contests its right and its cause, an enemy who would gladly rob it of its life and its hope. And this adversary, the Devil, yields before none but the righteous Judge; only God’s righteous judgment can strike down the Devil’s accusation and opposition and make it powerless. And often it seems as though God delays, as though it lasts so long, so dreadfully long, before He helps the poor soul in its distress.

Therefore the soul has such deep need to learn from the widow in our parable. It must do as she did; she plagued the unjust judge with her petitions so long until he helped her. Thus shall the Lord’s elect cry to Him day and night until He helps.

There are two things that we learn here about prayer. First, that it must be constant and unceasing. Second, that it must continue, even until it is fulfilled. It avails nothing to pray now and then, but what matters is to pray day and night. Neither does it avail to pray day and night for a time, for some weeks or months or years; but what matters is to pray without ceasing, until the answer comes, until the prayer is fulfilled.

Do you pray always, do you pray day and night? You do indeed have an adversary who watches constantly, an accuser who assails your right without ceasing. Do you pray just as unceasingly: Give me grace! Vindicate me? You do indeed know that you always need grace; you do indeed know that you always need the Lord’s righteousness. Or are you still so blinded that you think you can help yourself? But if you always need grace, and if it is only Christ’s righteousness that can strike down the Devil’s accusation, then pray, pray always, pray without ceasing, that Christ’s righteousness may be your righteousness. It stands written: "If God is for us, who can be against us? Who shall accuse God’s elect? God is He who justifies. Who is he that condemns? Christ is He who died, yea rather, who is also risen, who is also at God’s right hand, who also intercedes for us." Such promises are great and glorious, but how do they become yours, how can you appropriate them to yourself? They are appropriated to you through your filial Spirit of adoption, by whom you cry: Abba, Father! These promises are for the children, but the children are the praying souls in whom the Spirit works the unutterable groanings. If, then, prayer falls silent, and the sighs grow still, and the longing for Heaven vanishes, then the state of sonship is lost, and then the blessedness of the promise is gone. Faith is unceasing prayer, and poverty of spirit is an unceasing cry, and sonship is the blessed right to say: Abba, Father! Therefore pray without ceasing, pray day and night. Grace, grace, grace is the constant petition of believing prayer to God.

How long, you say, how long? Until you are heard, until you have received all that you need, until you have obtained vindication over your adversary; until the victory is won for ever, until the last enemy is overcome, until death is swallowed up for ever, until the crown is set upon your head, until the Son of Man comes. Then prayer ceases, when faith and hope are no more, but only the contemplation of eternal love.

So long as faith endures, prayer must endure. So long as the conflict continues, prayer must continue. So long as earthly life endures, it must be a life in prayer. For the victory is not won before the last enemy, death, has been struck down; and the inheritance, which your adversary would wrest from you, is not yours in full and outward possession before you have received the final answer from the Lord’s mouth: "Come, inherit the kingdom prepared for you before the foundation of the world was laid." Then you have received final judgment in your case; then your adversary is put to shame for ever and brought to silence.

Therefore endure unto the end and do not grow weary; only thus are obtained the crown of life and the wreath of righteousness.

But why then does the Savior say: "Nevertheless, when the Son of Man comes, will He find faith on the earth?" This is the serious word of warning to me and to you and to God’s Church. The times become more difficult for the elect, as it draws toward evening and the day declines. Prayer grows rarer because worldliness grows greater and men’s self-wisdom blinder. Already we hear the hoarse cries more and more loudly: "There is no God, there is no hearing of prayer, there is no salvation, there is no deliverance; all things happen according to a blind fate, according to an unchangeable predetermination, according to unchangeable laws of nature." It has come to pass according to the Lord’s word: "In the last days scoffers shall come, walking after their own lusts, and saying: Where is the promise of His coming? For since the fathers fell asleep, all things continue as they were from the beginning of the creation." It is this sluggish, cold evening fog that is spreading out over great portions of Christendom and dampening the fire of love and the warmth of prayer.

But then the Prophet says: "Upon your walls, O Jerusalem, I appoint watchmen; all the day and all the night, they shall never keep silence. You who remind the Lord, give yourselves no rest! And give Him no rest, until He establish, and until He make Jerusalem a praise in the earth."

O praying souls! It is you whom the Lord has appointed as watchmen; be faithful at your watch and be not silent day or night, never, until the Lord comes and His glory is revealed. Let not prayer grow still and faith die, but let your lamps be lit and your lights burning, lest all be darkness when the Son of Man comes!

